% 张三慧《大学物理学》第4版第25章《几何光学》原书思考题25.1--25.12
% V2主书PDF物理页339--340（原书322--323），和习题25.1--25.21区分。
\section*{第25章\quad 几何光学：思考题}
\paragraph{思考题 25.1}
在何种尺度与实验条件下可以忽略光的波动性，
把光传播近似视为沿直线传播？
\paragraph{思考题 25.2}
晴天树叶间小孔在地面形成圆形太阳斑；
日偏食时则出现许多月牙形光斑。
为什么斑的轮廓由太阳形状而不是小孔形状主导？
\paragraph{思考题 25.3}
照全身的平面穿衣镜，镜的最小竖直高度
是人身高的多少倍？是否与人离镜距离、眼睛高度有关？
\paragraph{思考题 25.4}
汽车后视镜与山路急弯处常用较大的凸面镜，
为什么比同等尺寸的平面镜有更宽视野？
镜中的后方车辆在视觉上通常比实际距离显得更远还是更近？
\paragraph{思考题 25.5}
草原炎热天气道路看似有水泊（下蜃景）
的物理原因是什么？海面远处的虚幻高像
（上蜃景）与它是否需要相同的空气折射率梯度？
\paragraph{思考题 25.6}
白光在三棱镜中因色散分出红光和紫光，
通常哪种光在玻璃中的相速率更大？
说明折射率与波长的关系。
\paragraph{思考题 25.7}
什么是光的全反射？在怎样的折射率与入射角条件下发生？
给出临界角公式，并讨论从低折射率侧射来的情形。
\paragraph{思考题 25.8}
把空气封入塑料薄膜制成饼状密封凸起，
浸没在折射率 \(n_w=1.33\) 的水中。
作为薄透镜，它使平行光会聚还是发散？
若两个球面曲率半径大小均为30cm，
估算薄透镜的焦距（空气相对于水的折射率）。
\paragraph{思考题 25.9}
列出真实物体通过理想薄凸透镜与薄凹透镜
形成像的各类位置、正倒、虚实及大小。
说明镜面成像是否适用相同符号约定。
\paragraph{思考题 25.10}
当物为位于透镜前的实物，
凸透镜在哪些像距范围内不能形成有限的实像？
凹透镜能形成实像或超出哪些虚像位置范围吗？
\paragraph{思考题 25.11}
通过球面镜或透镜实际观察时，
如何区分实像与虚像？它们能否分别用光屏承接？
\paragraph{思考题 25.12}
参照原图25.43(a)凸透镜及(b)凹透镜的光路，
若人眼想直接观察物体的像，
应处于出射光束的哪个区域？
作两条特征光线示意，并注明图中真实像/虚像位置依赖。
